\begin{lemma}
Let \(\eta>0\).  There is a threshold \(n_{\mathrm{base}}\) such that, for
every \(n>n_{\mathrm{base}}\), writing
\[
L=\log n,\quad
k_R=(1+\eta)\sqrt{nL},\quad
K=\noderef{TwoBiteNaturalIndependenceScale}(1+\eta,n),\quad k=(K:\mathbb R),
\]
\[
m_R=\frac{n}{L^2},\quad
M=\noderef{TwoBiteNaturalReducedVertexCount}(n),\quad m=(M:\mathbb R),
\]
\[
p=\frac12\sqrt{\frac Ln},\qquad
t_1=\frac{\sqrt{nL}}{\log L},
\]
we have
\[
0<m,\quad 0\le p\le\frac12,\quad 1\le2pm,
\]
\[
k_R\le k<k_R+1,\qquad m\le m_R<m+1,\qquad K\le n,\qquad k<2k_R,
\]
and
\[
0<\log L,\qquad
\frac{2k}{t_1}+1\le5(1+\eta)\log L.
\]
\end{lemma}
\begin{proof}
Choose \(n_{\mathrm{base}}\) so large that for every \(n>n_{\mathrm{base}}\),
\[
n>1,\quad L>1,\quad \log L>0,\quad k_R>2,\quad m_R>2,
\quad \sqrt{L/n}\le1,
\]
\[
\frac{k_R+1}{n}\le1,\qquad
\frac{\sqrt n}{2L^{3/2}}\ge1,\qquad \log L\ge1.
\]
These are eventual conditions because \(L\to\infty\), \(\log L\to\infty\),
\(\sqrt{nL}\to\infty\), \(n/L^2\to\infty\), \(\sqrt{L/n}\to0\),
\((k_R+1)/n\to0\), and \(\sqrt n/(2L^{3/2})\to\infty\).

Since \(L>1\), both \(k_R\) and \(m_R\) are nonnegative.  Applying
\(\noderef{NaturalParameterApproximation}\) with \(\kappa=1+\eta\) gives
\[
k_R\le k<k_R+1,\qquad m\le m_R<m+1.
\]
Because \(k_R>2\), we have \(k<k_R+1<2k_R\).  Because
\((k_R+1)/n\le1\), we also get \(k<n\), hence \(K\le n\).  Since
\(m_R>2\) and \(m_R<m+1\), we have
\[
\frac12m_R\le m\le m_R,\qquad 0<m.
\]
The definition of \(p\) and the inequality \(\sqrt{L/n}\le1\) give
\[
0\le p\le\frac12.
\]
Moreover,
\[
2pm\ge p m_R
=\frac12\sqrt{\frac Ln}\frac n{L^2}
=\frac{\sqrt n}{2L^{3/2}}\ge1.
\]
Finally, using \(k<2k_R=2(1+\eta)\sqrt{nL}\) and
\(t_1=\sqrt{nL}/\log L\),
\[
\frac{2k}{t_1}+1
<4(1+\eta)\log L+1.
\]
Since \(\log L\ge1\) and \(\eta>0\), the right side is at most
\[
5(1+\eta)\log L.
\]
This proves all claimed base-threshold estimates.
\end{proof}
